\chapter{الدوال}\label{ch:g10:functions}

\omterm{def:g10:functions:function}{الدالة} آلة تأخذ عددًا وتُخرج عددًا، وتعطي دائمًا الخرج نفسه من أجل
الدخل نفسه. والدوال هي الكائنات المركزية في التحليل الذي ستبنيه في
السنوات المقبلة (ابتداءً من \cref{ch:g11:quad})؛ ويضع هذا الفصل المفردات ---
\omterm{def:g10:functions:function}{مجموعة التعريف}، \omterm{def:g10:functions:function}{والصورة}، \omterm{def:g10:functions:graph}{والمنحنى}، والتغيرات --- ويمرّن على المهارة
الأساسية: قراءة المعلومات على \omterm{def:g10:functions:graph}{منحنى}.

\section{المفردات: الصور والسوابق}

\begin{definition}[الدالة]\label{def:g10:functions:function}
لتكن $D$ مجموعة من \omterm{def:g10:numbers:sets}{الأعداد الحقيقية}. \emph{الدالة}\index{دالة} $f$
المعرَّفة على $D$ ترفق بكل عدد $x \in D$ عددًا حقيقيًا واحدًا
بالضبط، يُكتب $f(x)$ ويُسمّى \emph{صورة}\index{صورة} العدد $x$.
والمجموعة $D$ هي \emph{مجموعة تعريف}\index{مجموعة التعريف} الدالة $f$. ونكتب
\[
f : x \longmapsto f(x).
\]
وإذا كان $f(x) = y$، فإن $x$ \emph{سابقة}\index{سابقة} للعدد $y$.
\end{definition}

\begin{remark}
لكل $x$ من \omterm{def:g10:functions:function}{مجموعة التعريف} \omterm{def:g10:functions:function}{صورة} واحدة بالضبط، لكن العدد $y$ قد تكون له
عدة سوابق أو لا تكون له \omterm{def:g10:functions:function}{سابقة}. فمن أجل $f(x) = x^2$: \omterm{def:g10:functions:function}{صورة} $3$ هي $9$، وللعدد
$9$ سابقتان هما $3$ و $-3$؛ أما العدد $-4$ فلا \omterm{def:g10:functions:function}{سابقة} له.
\end{remark}

\begin{example}[إيجاد مجموعة تعريف]\label{ex:g10:functions:domain}
عندما تُعطى \omterm{def:g10:functions:function}{دالة} بعبارة، تكون مجموعة تعريفها مجموعة الأعداد $x$ التي
تكون العبارة من أجلها ذات معنى.
\begin{itemize}
  \item $f(x) = 3x^2 - 5x + 1$ لها معنى من أجل كل $x$: \omterm{def:g10:functions:function}{مجموعة التعريف} هي
        $\R$.
  \item $g(x) = \dfrac{1}{x - 2}$ تستلزم $x - 2 \neq 0$: \omterm{def:g10:functions:function}{مجموعة التعريف} هي
        كل \omterm{def:g10:numbers:sets}{الأعداد الحقيقية} عدا $2$.
  \item $h(x) = \sqrt{x - 3}$ تستلزم $x - 3 \geq 0$: \omterm{def:g10:functions:function}{مجموعة التعريف} هي
        $\intco{3}{+\infty}$.
\end{itemize}
\end{example}

\begin{example}[حساب الصور والسوابق]\label{ex:g10:functions:images}
لتكن $f(x) = x^2 - 4x + 3$، معرَّفة على $\R$.

\emph{\omterm{def:g10:functions:function}{صورة} $5$:} نعوّض $x = 5$:
$f(5) = 25 - 20 + 3 = 8$.

\emph{سوابق $3$:} نحل $f(x) = 3$:
\[
x^2 - 4x + 3 = 3
\quad\Longleftrightarrow\quad
x^2 - 4x = 0
\quad\Longleftrightarrow\quad
x(x - 4) = 0,
\]
إذن سابقتا $3$ هما $0$ و $4$.
\end{example}

\section{منحنى دالة}

\begin{definition}[المنحنى]\label{def:g10:functions:graph}
في معلم، \emph{منحنى}\index{منحنى} \omterm{def:g10:functions:function}{الدالة} $f$ (أو \emph{تمثيلها البياني}) هو
مجموعة كل النقط $(x, f(x))$ من أجل $x$ في \omterm{def:g10:functions:function}{مجموعة التعريف}. وبعبارة
أخرى، تقع النقطة $(x, y)$ على المنحنى بالضبط عندما $y = f(x)$.
\end{definition}

\begin{omfigure}
\begin{tikzpicture}
\begin{axis}[omaxis,
  width=8.6cm, height=6.0cm,
  xmin=-2.6, xmax=4.6, ymin=-3.4, ymax=4.4,
  xtick={-2,-1,1,2,3,4}, ytick={-3,-2,-1,1,2,3,4}]
  \addplot[omDef, thick, domain=-2.2:4.2, samples=120]
    {0.25*x^3 - 0.75*x^2 - x + 2};
  \draw[dashed, omThm] (axis cs:3,0) -- (axis cs:3,1.25)
    -- (axis cs:0,1.25);
  \fill[omThm] (axis cs:3,1.25) circle (1.5pt);
  \node[omThm, font=\small, above right] at (axis cs:3,1.2)
    {$(3,\,f(3))$};
\end{axis}
\end{tikzpicture}

{\small قراءة \omterm{def:g10:functions:function}{صورة} على \omterm{def:g10:functions:graph}{منحنى}: انطلق من $x = 3$ على المحور
الأفقي، وارتفع شاقوليًا حتى \omterm{def:g10:functions:graph}{المنحنى}، ثم امضِ أفقيًا حتى المحور الشاقولي
لتقرأ $f(3) = 1.25$.}
\end{omfigure}

\begin{method}[قراءة منحنى]\label{met:g10:functions:reading}
على \omterm{def:g10:functions:graph}{منحنى} \omterm{def:g10:functions:function}{الدالة} $f$:
\begin{enumerate}
  \item \emph{\omterm{def:g10:functions:function}{صورة} $a$}: امضِ شاقوليًا من $x = a$ على المحور الأفقي
        حتى \omterm{def:g10:functions:graph}{المنحنى}، ثم أفقيًا حتى المحور الشاقولي؛ والقيمة
        المقروءة هناك هي $f(a)$؛
  \item \emph{سوابق $b$}: ارسم المستقيم الأفقي $y = b$؛
        والسوابق هي الفواصل $x$ لكل نقط تقاطعه
        مع \omterm{def:g10:functions:graph}{المنحنى}؛
  \item \emph{حلول $f(x) = k$}: هي نفسها سوابق $k$؛
  \item \emph{حلول $f(x) \leq k$}: هي الأعداد $x$ التي يكون \omterm{def:g10:functions:graph}{المنحنى} عندها
        على المستقيم $y = k$ أو تحته.
\end{enumerate}
\end{method}

\begin{omfigure}
\begin{tikzpicture}
\begin{axis}[omaxis,
  width=8.6cm, height=6.0cm,
  xmin=-3.4, xmax=3.4, ymin=-2.4, ymax=4.4,
  xtick={-2,2}, xticklabels={$x_1$,$x_2$}, ytick={2}, yticklabels={$k$}]
  \addplot[omDef, thick, domain=-3:3, samples=120] {0.5*x^2};
  \addplot[omThm, thick, domain=-3.2:3.2] {2};
  \fill (axis cs:-2,2) circle (1.5pt);
  \fill (axis cs:2,2) circle (1.5pt);
  \draw[dashed] (axis cs:-2,2) -- (axis cs:-2,0);
  \draw[dashed] (axis cs:2,2) -- (axis cs:2,0);
\end{axis}
\end{tikzpicture}

{\small حل $f(x) = k$ بيانيًا: الحلان $x_1$ و $x_2$ هما
فاصلا نقطتي \omterm{def:g10:numbers:interunion}{تقاطع} \omterm{def:g10:functions:graph}{المنحنى} مع المستقيم الأفقي
$y = k$. وهنا تتحقق $f(x) \leq k$ على $\intcc{x_1}{x_2}$، حيث يكون
\omterm{def:g10:functions:graph}{المنحنى} تحت المستقيم.}
\end{omfigure}

\begin{example}\label{ex:g10:functions:membership}
هل تنتمي النقطة $A(2, 5)$ إلى \omterm{def:g10:functions:graph}{منحنى} \omterm{def:g10:functions:function}{الدالة} $f(x) = x^2 + 1$؟ احسب
$f(2) = 4 + 1 = 5$: نعم، لأن $f(2)$ يساوي الترتيب $y$ للنقطة $A$.
أما النقطة $B(3, 8)$ فلا تنتمي إلى \omterm{def:g10:functions:graph}{المنحنى}، لأن $f(3) = 10 \neq 8$.
\end{example}

\section{تغيرات دالة}

\begin{definition}[متزايدة، متناقصة]\label{def:g10:functions:variations}
لتكن $f$ معرَّفة على \omterm{def:g10:numbers:interval}{فترة} $I$.
\begin{itemize}
  \item تكون $f$ \emph{متزايدة}\index{دالة متزايدة} على $I$ إذا كان،
        من أجل كل $u, v \in I$، $u < v$ يستلزم $f(u) \leq f(v)$: أي أن
        الخرج ينمو مع الدخل، ويصعد \omterm{def:g10:functions:graph}{المنحنى} من اليسار إلى
        اليمين.
  \item وتكون $f$ \emph{متناقصة}\index{دالة متناقصة} على $I$ إذا كان
        $u < v$ يستلزم $f(u) \geq f(v)$: أي أن \omterm{def:g10:functions:graph}{المنحنى} يهبط من اليسار إلى
        اليمين.
\end{itemize}
ومع المتراجحات التامة $f(u) < f(v)$ (على الترتيب $>$) نقول إنها
متزايدة \emph{تمامًا} (على الترتيب متناقصة تمامًا).
\end{definition}

\begin{definition}[جدول التغيرات]\label{def:g10:functions:table}
\emph{جدول التغيرات}\index{جدول التغيرات} يلخّص الفترات التي تتزايد
عليها $f$ والتي تتناقص عليها، بسهمين $\nearrow$ و
$\searrow$، ويسجل قيم $f$ عند نقط الانعطاف.
\end{definition}

\begin{example}\label{ex:g10:functions:table}
لننظر في \omterm{def:g10:functions:function}{الدالة} الممثَّلة أدناه على $\intcc{-2}{4}$.

\begin{omfigure}
\begin{tikzpicture}
\begin{axis}[omaxis,
  width=8.6cm, height=5.4cm,
  xmin=-2.8, xmax=4.6, ymin=-2.4, ymax=3.8,
  xtick={-2,1,4}, ytick={-1.5,3}, yticklabels={$-1.5$,$3$}]
  \addplot[omDef, thick, domain=-2:4, samples=120] {-0.5*(x-1)^2 + 3};
  \fill (axis cs:-2,-1.5) circle (1.5pt);
  \fill (axis cs:4,-1.5) circle (1.5pt);
  \fill (axis cs:1,3) circle (1.5pt);
\end{axis}
\end{tikzpicture}

{\small \omterm{def:g10:functions:function}{دالة} معرَّفة على $\intcc{-2}{4}$، \omterm{def:g10:functions:variations}{متزايدة} ثم \omterm{def:g10:functions:variations}{متناقصة}.}
\end{omfigure}

بقراءة \omterm{def:g10:functions:graph}{المنحنى}: تتزايد $f$ من $f(-2) = -1.5$ حتى قيمتها العظمى
$f(1) = 3$، ثم تتناقص حتى $f(4) = -1.5$. وجدول تغيراتها هو
\begin{center}
\renewcommand{\arraystretch}{1.3}
\begin{tabular}{c|ccccc}
$x$ & $-2$ & & $1$ & & $4$ \\
\hline
$f$ & $-1.5$ & $\nearrow$ & $3$ & $\searrow$ & $-1.5$ \\
\end{tabular}
\end{center}
\end{example}

\begin{example}[البرهان على تغير]\label{ex:g10:functions:proof}
بيّن أن $f(x) = 3x + 1$ \omterm{def:g10:functions:variations}{متزايدة} تمامًا على $\R$. خذ $u < v$
كيفما كانا وقارن الصورتين:
\[
f(v) - f(u) = (3v + 1) - (3u + 1) = 3(v - u) > 0,
\]
لأن $v - u > 0$. إذن $f(u) < f(v)$: \omterm{def:g10:functions:variations}{فالدالة متزايدة}
تمامًا. ويعطي الحساب نفسه، مع ميل سالب، مثلًا
$g(x) = -2x + 5$، ما يلي: $g(v) - g(u) = -2(v-u) < 0$: \omterm{def:g10:functions:variations}{متناقصة} تمامًا.
\end{example}

\section{القيم الحدية}

\begin{definition}[القيمة العظمى، القيمة الصغرى]\label{def:g10:functions:extrema}
لتكن $f$ معرَّفة على $D$ وليكن $a \in D$. تكون القيمة $f(a)$
\emph{القيمة العظمى}\index{قيمة عظمى} \omterm{def:g10:functions:function}{للدالة} $f$ على $D$ إذا كان $f(x) \leq f(a)$ من أجل كل
$x \in D$؛ وتكون \emph{القيمة الصغرى}\index{قيمة صغرى} إذا كان
$f(x) \geq f(a)$ من أجل كل $x \in D$. وعلى \omterm{def:g10:functions:graph}{المنحنى}، هما أعلى نقطة
وأخفض نقطة فيه.
\end{definition}

\begin{example}\label{ex:g10:functions:extremum}
بيّن أن $f(x) = (x - 1)^2 + 2$ قيمتها الصغرى $2$ على $\R$، وتبلغها عند
$x = 1$. من أجل كل $x$، يكون المربع $(x-1)^2$ $\geq 0$، إذن
$f(x) = (x-1)^2 + 2 \geq 2$؛ و$f(1) = 0 + 2 = 2$، فالقيمة $2$
مبلوغة فعلًا. وهاتان الواقعتان معًا تبرهنان على أن $2$ هي \omterm{def:g10:functions:extrema}{القيمة الصغرى}.
\end{example}

\begin{omfigure}
\begin{tikzpicture}
\begin{axis}[omaxis,
  width=7.6cm, height=5.2cm,
  xmin=-1.8, xmax=3.8, ymin=-0.9, ymax=6.4,
  xtick={1}, ytick={2}]
  \addplot[omDef, thick, domain=-1.0:3.0, samples=100] {(x-1)^2 + 2};
  \addplot[omThm, dashed, domain=-1.6:3.6] {2};
  \fill (axis cs:1,2) circle (1.5pt);
  \node[below right, font=\small] at (axis cs:1,2) {$(1,2)$};
\end{axis}
\end{tikzpicture}

{\small لا ينزل \omterm{def:g10:functions:graph}{منحنى} $f(x) = (x-1)^2 + 2$ أبدًا تحت المستقيم المتقطع
$y = 2$، ويلامسه عند $x = 1$: إذن \omterm{def:g10:functions:extrema}{القيمة الصغرى} \omterm{def:g10:functions:function}{للدالة} $f$ هي $2$.}
\end{omfigure}

\section{تمارين}

\begin{exercise}[$\star$]\label{exo:g10:functions:1}
لتكن $f(x) = 2x^2 - 3x + 1$. احسب صور الأعداد $0$ و $2$ و $-1$ و
$\frac12$.
\end{exercise}

\begin{exercise}[$\star$]\label{exo:g10:functions:2}
أعطِ \omterm{def:g10:functions:function}{مجموعة تعريف} كل \omterm{def:g10:functions:function}{دالة}:
\[
f(x) = 5x - 2, \qquad
g(x) = \frac{3}{x + 4}, \qquad
h(x) = \sqrt{2x - 6}, \qquad
k(x) = \frac{1}{x^2 + 1}.
\]
\end{exercise}

\begin{exercise}[$\star$]\label{exo:g10:functions:3}
لتكن $f(x) = x^2 - 2x$. أوجد كل سوابق $0$، ثم $3$، ثم $-1$.
\end{exercise}

\begin{exercise}[$\star$]\label{exo:g10:functions:4}
هل تنتمي النقطة $A(-1, 4)$ إلى \omterm{def:g10:functions:graph}{منحنى} \omterm{def:g10:functions:function}{الدالة} $f(x) = x^2 - 2x + 1$؟
وماذا عن النقطة $B(2, 1)$؟ برّر بحساب.
\end{exercise}

\begin{exercise}[$\star$]\label{exo:g10:functions:5}
\omterm{def:g10:functions:function}{لدالة} $g$ معرَّفة على $\intcc{-3}{5}$ \omterm{def:g10:functions:table}{جدول التغيرات}
\begin{center}
\renewcommand{\arraystretch}{1.3}
\begin{tabular}{c|ccccccc}
$x$ & $-3$ & & $0$ & & $3$ & & $5$ \\
\hline
$g$ & $1$ & $\searrow$ & $-2$ & $\nearrow$ & $4$ & $\searrow$ & $0$ \\
\end{tabular}
\end{center}
\begin{enumerate}
  \item ما \omterm{def:g10:functions:extrema}{القيمة العظمى} \omterm{def:g10:functions:extrema}{والقيمة الصغرى} \omterm{def:g10:functions:function}{للدالة} $g$ على $\intcc{-3}{5}$؟
  \item قارن $g(-1)$ و $g(-0.5)$ دون معرفة قيمتيهما.
  \item كم حلًّا على الأكثر \omterm{def:g10:algebra:equation}{للمعادلة} $g(x) = 0$ على
        كل \omterm{def:g10:numbers:interval}{فترة} من فترات الجدول؟
\end{enumerate}
\end{exercise}

\begin{exercise}[$\star\star$]\label{exo:g10:functions:6}
لتكن $f(x) = -2x + 7$. بيّن، بمقارنة $f(u)$ و $f(v)$ من أجل $u < v$،
أن $f$ \omterm{def:g10:functions:variations}{متناقصة} تمامًا على $\R$.
\end{exercise}

\begin{exercise}[$\star\star$]\label{exo:g10:functions:7}
بيّن أن \omterm{def:g10:functions:function}{للدالة} $f(x) = x^2 + 6x + 11$ \omterm{def:g10:functions:extrema}{قيمة صغرى} على $\R$، وأعطِ
قيمتها والنقطة التي تُبلغ عندها. (إرشاد: اكتب
$f(x) = (x + 3)^2 + c$ من أجل الثابت $c$ المناسب.)
\end{exercise}

\begin{exercise}[$\star\star$]\label{exo:g10:functions:8}
حديقة مستطيلة محيطها $40$ م. ليكن $x$ عرضها بالمتر.
\begin{enumerate}
  \item عبّر عن الطول، ثم عن المساحة $A(x)$، بدلالة $x$.
        من أجل أي قيم $x$ يكون لهذا معنى؟
  \item احسب $A(4)$ و $A(8)$ و $A(10)$ و $A(12)$ و $A(16)$. ماذا
        تخمّن عن الشكل ذي المساحة الأكبر؟
\end{enumerate}
\end{exercise}

\begin{exercise}[$\star\star$]\label{exo:g10:functions:9}
لتكن $f(x) = \dfrac{1}{x^2 + 1}$، معرَّفة على $\R$.
\begin{enumerate}
  \item بيّن أن $0 < f(x) \leq 1$ من أجل كل \omterm{def:g10:numbers:sets}{عدد حقيقي} $x$.
  \item هل \omterm{def:g10:functions:function}{للدالة} $f$ \omterm{def:g10:functions:extrema}{قيمة عظمى} على $\R$؟ \omterm{def:g10:functions:extrema}{وقيمة صغرى}؟ برّر.
\end{enumerate}
\end{exercise}

\begin{exercise}[$\star\star\star$]\label{exo:g10:functions:10}
لتكن $f(x) = x^2$ على $\R$.
\begin{enumerate}
  \item ليكن $0 \leq u < v$. حلّل $f(v) - f(u)$ إلى عوامل واستنتج أن $f$
        \omterm{def:g10:functions:variations}{متزايدة} تمامًا على $\intco{0}{+\infty}$.
  \item كيّف الحجة لتبيّن أن $f$ \omterm{def:g10:functions:variations}{متناقصة} تمامًا على
        $\intoc{-\infty}{0}$.
\end{enumerate}
\end{exercise}

\section{مسألة: آلات تُطعَم من خرجها}

\begin{problem}[{مسألة نهاية الأسبوع --- النقط الثابتة، والتكرار،
وثلاث آلات شهيرة: محسّنة هيرون، وآلة التربيع، وآلة
البَرَد التي لم تُحلّ}]
\label{pb:g10:functions:1}
\omterm{def:g10:functions:function}{الدالة} آلة: يدخل عدد ويخرج عدد.
وأمتع التجارب هي التي تُطعم الآلة \emph{خرجها
هي}، مرة بعد مرة --- وعندئذٍ يسيطر سؤالان:
هل يستقر المسار عند موضع ما (\emph{\omterm{pb:g10:functions:1}{نقطة ثابتة}})،
وكيف يسوقه تغير الآلة إلى هناك؟ إحدى آلات هذه
المسألة تصقل الجذور التربيعية منذ ألفي
سنة؛ وأخرى تحرس أشهر مسألة لم تُحلّ
في الرياضيات الابتدائية.

\medskip
\noindent\textbf{الجزء الأول --- آلة هيرون.}
لننظر في \omterm{def:g10:functions:function}{الدالة}
\[
f(x) = \frac12\left(x + \frac2x\right).
\]
\begin{enumerate}
  \item أعطِ \omterm{def:g10:functions:function}{مجموعة تعريف} $f$، واحسب $f(1)$ و $f(2)$ و
        $f\!\left(\frac32\right)$ على \omterm{def:g10:functions:function}{صورة} كسور مضبوطة.
  \item أوجد كل سوابق $\frac32$: حل
        $f(x) = \frac32$ (تخلّص من المقام واستعمل
        طرائق \cref{pb:g10:algebra:1}).
  \item \emph{النقطة الثابتة}\index{نقطة ثابتة} \omterm{def:g10:functions:function}{للدالة} $f$ هي
        عدد لا تغيّره الآلة: $f(x) = x$. أوجد النقطتين
        الثابتتين \omterm{def:g10:functions:function}{للدالة} $f$. ومَن ذلك المعرفة القديمة التي تحرس
        الموجبة منهما؟
  \item انطلاقًا من $x = 1$، أطعم الآلة خرجها
        ثلاث مرات وسجّل النتائج المضبوطة. أي متتالية
        من مسألة نهاية الأسبوع عن الصمم في الكتاب السابق أعدتَ بناءها --- وما كانت
        ``وصفة هيرون'' بلغة الدوال؟
  \item في رسم واحد، اُرسم \omterm{def:g10:functions:graph}{منحنى} $f$ (من أجل $x > 0$)
        والمستقيم القطري $y = x$. أين تقع النقطتان
        الثابتتان في هذا الرسم؟ (قارن
        بمسألة نهاية الأسبوع عن الميزانين الحراريين في الكتاب السابق: فقراءة درجة الحرارة نفسها
        في السلّمين كانت الفكرة نفسها.)
\end{enumerate}

\medskip
\noindent\textbf{الجزء الثاني --- لماذا لا تخطئ الآلة.}
\begin{enumerate}[resume]
  \item من أجل $0 < u < v$، بيّن بتوحيد المقامات أن
        \[
        f(v) - f(u) = \frac{(v - u)(uv - 2)}{2uv} ,
        \]
        واستنتج تغيرات $f$ على
        $\intoo{0}{+\infty}$: \omterm{def:g10:functions:variations}{متناقصة} حتى نقطة، ثم
        \omterm{def:g10:functions:variations}{متزايدة} --- فما تلك النقطة؟ (التقنية نفسها المستعملة في
        \cref{exo:g10:functions:10}.)
  \item استنتج أن \omterm{def:g10:functions:function}{للدالة} $f$ على $\intoo{0}{+\infty}$
        \omterm{def:g10:functions:extrema}{قيمة صغرى} تساوي $\sqrt2$ بالضبط، وتُبلغ عند
        $\sqrt2$. ماذا يقول ذلك عن \emph{كل} خرج
        من الآلة (من أجل دخل موجب)؟
  \item بيّن أنه إذا كان $x > \sqrt2$ فإن
        $\sqrt2 < f(x) < x$: أي أن أي تخمين مفرط في الكبر يتحسّن دائمًا،
        ولا يُتجاوز أبدًا. (من أجل $f(x) < x$، قارن $\frac2x$
        بالعدد $x$.)
  \item اجمع الأسئلة 6--8 في جدول تغيرات $f$
        على $\intoo{0}{+\infty}$
        (\cref{def:g10:functions:table})، مع \omterm{def:g10:functions:extrema}{القيمة الصغرى}.
  \item آلة ثانية: $g(x) = x^2$. نقطتاها الثابتتان هما
        $0$ و $1$ (تحقق). كرّر $g$ أربع مرات انطلاقًا من
        $x = 0.9$، ثم أربع مرات انطلاقًا من $x = 1.1$ (بالحاسبة،
        بثلاثة أرقام عشرية). إحدى النقطتين الثابتتين جاذبة والأخرى
        طاردة: فأيهما أيهما؟
\end{enumerate}

\medskip
\noindent\textbf{الجزء الثالث --- آلة البَرَد.}
على \omterm{def:g10:numbers:sets}{الأعداد الصحيحة} الطبيعية، نعرّف: $h(n) = \frac n2$ إذا كان $n$ زوجيًا،
و $h(n) = 3n + 1$ إذا كان $n$ فرديًا. وترتد الأعداد تحت $h$ مثل
حبّات البَرَد في سحابة عاصفة.
\begin{enumerate}[resume]
  \item احسب رحلة $7$ الكاملة تحت $h$، نزولًا إلى $1$.
        كم خطوة تستغرق، وإلى أي علوّ تحلّق؟
  \item ابدأ رحلة $27$ واحسب عشر خطوات. (تدوم
        رحلته الكاملة $111$ خطوة وتبلغ ذروتها عند $9\,232$ ---
        انطلاقًا من $27$.) أي درس عن القواعد
        البسيطة يعلّمنا إياه العددان $7$ و $27$؟
  \item فسّر لماذا تهوي إلى $1$ مباشرةً كل رحلة تبلغ قوة
        للعدد $2$، وماذا يحدث عند $1$
        (احسب $h(1)$ و $h(4)$ و $h(2)$). هل \omterm{def:g10:functions:function}{للدالة} $h$
        \omterm{pb:g10:functions:1}{نقطة ثابتة} بين \omterm{def:g10:numbers:sets}{الأعداد الصحيحة} الموجبة؟ وماذا
        تشكّل نهاية كل رحلة مرصودة بدلًا من ذلك؟
  \item تدّعي \emph{حدسية كولاتز} أن كل عدد
        بداية يبلغ $1$ في النهاية. وقد تحققت الحواسيب من ذلك
        إلى ما هو أبعد بكثير من $10^{20}$؛ لكن ما من إنسان برهن عليها. أي
        نوعين من الاكتشاف \emph{يدحضانها}؟ ولماذا
        لا يغلق التحقق الهائل الملف
        (لقد قرعت مسألة نهاية الأسبوع عن عرّاف أعواد الثقاب في الكتاب السابق هذا الجرس)؟
  \item ``لم تنضج الرياضيات بعد لأسئلة كهذه''، قال
        بول إردوش عن هذه الحدسية. في جملة واحدة: ما الذي
        يفصل آلة البَرَد عن آلة هيرون، حيث حسمت
        الأسئلة 6--8 كل شيء؟
\end{enumerate}

\medskip
\noindent\textbf{الجزء الرابع --- دليل استعمال الآلات.}
\begin{enumerate}[resume]
  \item مجموعات التعريف على \omterm{def:g10:functions:function}{صورة} فترات (بلغة
        \cref{pb:g10:numbers:1}): أعطِ مجموعات تعريف
        $x \mapsto \sqrt{x - 3}$، \
        $x \mapsto \frac{1}{x^2 - 9}$، \
        $x \mapsto \sqrt{9 - x^2}$.
  \item عودًا إلى \omterm{def:g10:functions:function}{دالة} هيرون $f$: أوجد سوابق $3$ المضبوطة
        (حل $f(x) = 3$ بإتمام المربع).
  \item يُقذف حجر إلى أعلى: ارتفاعه
        $H(t) = 20t - 5t^2$ مترًا بعد $t$ ثانية. على أي
        \omterm{def:g10:numbers:interval}{فترة} زمنية يكون للعبارة معنى فيزيائي؟
        أتمّ المربع لإيجاد الارتفاع الأعظمي ولحظته
        (\cref{def:g10:functions:extrema})، وأعطِ
        \omterm{def:g10:functions:table}{جدول التغيرات}.
  \item كم من الوقت يبقى الحجر على ارتفاع $15$ م على الأقل؟ (حل
        $H(t) \geq 15$ بجدول إشارة،
        \cref{met:g10:algebra:signtable}.)
  \item الخاتمة: شغّلت هذه المسألة أربع آلات --- محسّنة
        هيرون، وآلة التربيع، وآلة البَرَد، وارتفاع
        الحجر. في جملتين أو ثلاث، بيّن ما تملكه كل
        \omterm{def:g10:functions:function}{دالة} (\omterm{def:g10:functions:function}{مجموعة تعريف}، وقاعدة، \omterm{def:g10:functions:graph}{ومنحنى}) والسؤالين
        الكبيرين اللذين طرحتهما هذه المسألة على كل منها (أين يستقر
        التكرار؟ وكيف تتغير \omterm{def:g10:functions:function}{الدالة}؟) ---
        وهما سؤالان تتناولهما الفصول المقبلة واحدًا واحدًا.
\end{enumerate}
\end{problem}
